
\documentclass[../../../ps_q.tex]{subfiles}

\begin{document}
    真空中に$xy$平面を定め，$y$軸上の点$(0,\myspacea L)$に正の点電荷A（電荷$Q$）を固定し，%
    さらに$y$軸上の点$(0,\myspacea -L)$に負の点電荷B（電荷$-Q$）を固定した．%
    クーロンの法則の比例定数を$k_0$とし，クーロン力以外の力の影響は考えない．%
    点$(x,\myspacea y)$における電場の$x$成分を$E_x(x,\myspacea y)$，%
    $y$成分を$E_y(x,\myspacea y)$，%
    電位（無限遠基準）を$\phi(x,\myspacea y)$と表す．

    \begin{enumerate}
        \item $E_x(L,\myspacea L)$，$E_y(L,\myspacea L)$を求めよ．
        \item $E_x(x,\myspacea 0)$，$E_y(x,\myspacea 0)$，%
            $\phi(x,\myspacea 0)$を求めよ．
        \item $E_x(0,\myspacea y)$，$E_y(0,\myspacea y)$，%
            $\phi(0,\myspacea y)$を求めよ．%
            ただし，$y \neq \pm L$とする．
        \item $y$軸に沿ってなめらかに動ける正の点電荷C（質量$m$，電荷$q$）を，%
            点$(0,\myspacea 2L)$で静かに放したところ，Cは無限遠方へ飛び去った．%
            Cが無限遠に到達したときの速さ$v_\infty$を求めよ．
    \end{enumerate}
\end{document}


